\documentclass[11pt]{article}
\usepackage[spanish, es-tabla, es-lcroman]{babel}

%Paquetes básicos
\usepackage{multicol}
\usepackage{physics} 
\usepackage{float}
\usepackage{gensymb}
\usepackage{siunitx} 
\usepackage{enumerate} 
\usepackage{url}
\usepackage{pgfplots}
\usepackage{pgfplotstable}
\usepackage{tikz,pgfplots}
\usepackage{amsmath}  
\usepackage{wasysym} 
\usepackage{geometry}
\usepackage{mdframed}
%Fuente: Helvetica
\usepackage[scaled]{helvet}
\usepackage[T1]{fontenc}
\renewcommand\familydefault{\sfdefault}
\usepackage[eulergreek]{sansmath}
\usepackage[frenchmath]{newtxsf}
\renewcommand{\familydefault}{\sfdefault}
\usepackage{mathastext}
\usepackage{lipsum}
\usepackage{apacite}
\usepackage{natbib}

%Paquetes necesarios para la elaboración de Figuras con el ambiente tikzpicture
\usepackage{tikz} 
\tikzset{font=\fontfamily{phv}\selectfont}
\usepackage{amsmath,tikz}
\usepackage{tikz-3dplot}
\usetikzlibrary{calc,positioning,intersections}
\usetikzlibrary{positioning,shapes.misc}
\tdplotsetmaincoords{80}{120}
\usetikzlibrary{decorations.markings} 
\usetikzlibrary{decorations.pathmorphing}
\usetikzlibrary{shapes, shapes.geometric}
\usetikzlibrary{mindmap,trees}
\usetikzlibrary{backgrounds, fit, positioning}
\tikzstyle{flecha} = [thick,->,>=stealth]


%Geometría de la página y encabezados
\geometry{a4paper,top=3cm,bottom=2.5cm,right=2.5cm,left=2.5cm}
\pgfplotsset{compat=1.14}
\pgfplotsset{/pgf/number format/use comma}
%Estilo de encabezados
\usepackage{fancyhdr}
\pagestyle{fancy}
\fancyhf{}
\lhead{LAB-IA-I}
\rhead{FCQ-UNA}
\cfoot{\thepage}




%-----------------------
%AQUÍ SE DEBE REEMPLAZAR EL SÍMBOLO NUMERAL POR EL NÚMERO DE LA PRÁCTICA
\chead{Práctica N$^{\circ}$\#}
%-----------------------



%Configuración del estilo de los títulos de las secciones y subsecciones
\usepackage{titlesec}
\titleformat{\section}{\normalfont\normalsize\bfseries\centering}{\thesection.}{0.5em}{}
\titleformat{\subsection}{\normalfont\normalsize\bfseries}{\thesubsection.}{0.5em}{}



%-----------------------
%AQUÍ SE DEBE REEMPLAZAR LO ESCRITO POR EL TÍTULO DE LA PRÁCTICA
\title{\LARGE\textbf{Título de la práctica}}
%-----------------------



%-----------------------
%Y LOS NOMBRES DE CADA UNO DE LOS INTEGRANTES DEL GRUPO
\author{\normalsize{Apellido, Nombre $\cdot$ Apellido, Nombre $\cdot$ Apellido, Nombre $\cdot$ Apellido, Nombre}}
%-----------------------


\date{\small{\today}}

\begin{document}
\renewcommand{\BOthers}[1]{et al.\hbox{}}


\maketitle


\hrule

\begin{abstract}
\noindent %En general, el resumen no contiene sangría, el comando noident se encarga de no poner la sangría al párrafo
El resumen de la práctica no debe contener más de 250 palabras, debe ser conciso y brindar una idea clara sobre el trabajo y sus conclusiones. Además, debe contener las palabras claves del trabajo.

\noindent\textit{\textbf{Palabras claves:} palabra clave 1, palabra clave 2, palabra clave 3.}

\end{abstract}

\hrule


\begin{multicols}{2}
\section{Introducción}
En esta sección se desarrolla el contexto del trabajo, incluye una exhaustiva revisión bibliográfica, se exponen los objetivos del trabajo y siempre se incluyen las referencias bibliográficas. 

Como ejemplo, se puede mencionar la ecuación de flujo incompresible (Ecuación \ref{ec:1}), más bien conocida como ecuación de Bernoulli \citep{streeter}:
\begin{equation}\label{ec:1}
\dfrac{V^{2}\cdot \rho}{2}+P+\rho\cdot g \cdot z=constante
\end{equation}

Donde, $V$ es la velocidad del fluido, $\rho$ es la densidad del fluido, $P$ es la presión a lo largo de la línea de corriente, $g$ es la aceleración de la gravedad y $z$ es la altura en la cota de referencia.
\section{Metodología}

En esta sección se debe detallar el procedimiento ejecutado durante el trabajo y la técnica empleada debe ser debidamente referenciada. Los materiales y equipos utilizados durante la experiencia deben ser mencionados. 

\section{Resultados y discusiones}
Los resultados obtenidos deben ser presentados de manera ordenada, clara y concisa, preferentemente, en el mismo orden en el cual se planteó la metodología. 

La presentación de los resultados se puede realizar en forma de tablas, como ejemplo se observa la Tabla \ref{tab:tab1}.

\begin{table}[H]
    \centering
    \caption{Coeficientes representativos de pérdidas de cabeza ($K$), para diferentes accesorios}
\begin{tabular}{lr}
\hline
\multicolumn{1}{c}{Accesorio} & \multicolumn{1}{c}{$K$}\\
\hline
Codo en U & 2,20\\
Codo estándar & 0,90\\
Codo de radio medio & 0,75\\
Codo de radio largo & 0,60\\
\hline
\end{tabular}
    \label{tab:tab1}
    \begin{minipage}{\columnwidth}
    \centering Fuente: \citet{streeter}.
    \end{minipage}
\end{table}

Otra forma de presentar los resultados es por medio de gráficos, los mismos deben elaborarse claramente, los ejes deben contener nombres y las correspondientes unidades de medida, como se observa en la Figura \ref{fig:1}.
\end{multicols}

\begin{figure}[H]
    \centering
    \caption{Relación entre los caudales reales y experimentales durante la descarga de agua a través de un vertedero de geometría triangular}
\begin{mdframed}
\centering 
    \begin{tikzpicture}[font=\sffamily\sansmath\footnotesize]
\begin{axis}[
    xmin=0, ymin=0, xmax=0.25, ymax=0.15,
	xlabel=Caudal experimental (L/s),
	ylabel=Caudal real (L/s),
	width=0.6\columnwidth,height=7cm,
    ]
\addplot[color=blue,mark=x,smooth] coordinates {
(0.02,0.01)
(0.10,0.06)
(0.12,0.07)
(0.17,0.10)
(0.23,0.13)
};
\end{axis}
\end{tikzpicture}
\end{mdframed}
    \label{fig:1}
\begin{minipage}{\columnwidth}
\centering Fuente: Elaboración propia.
\end{minipage}
\end{figure}
\begin{multicols}{2}


Expuestos los resultados, se debe presentar un análisis completo y profundo de los resultados, coherente y fundamentado en bases teóricas, no deben dejarse puntos sin analizar o con explicaciones ambiguas.

\section{Conclusiones}

La conclusión se realiza en base a los resultados obtenidos y los objetivos trazados para la práctica, debe estar redactada de una forma clara, correcta y concisa.

%La ventaja de utilizar la plataforma Overleaf para la redacción de los informes de laboratorio e inclusive cualquier otro documento técnico, es la posibilidad de auto-redacción de las referencias bibliográficas, mediante el empleo de un archivo tipo .bib encontrado en el panel de control junto con el archivo con el código principal, basta con seleccionar el estilo de citación y el archivo que contiene las citaciones para que Overleaf elabore la bibliografía.

\bibliographystyle{apacite} 
\bibliography{referencias.bib}

\end{multicols}



\end{document}
